% DOCUMENT FORMAT ======================= -*- mode: LaTeX; coding: utf-8 -*- ===

\documentclass[diploma]{softlab-thesis}


% PACKAGE SETTINGS =============================================================

\usepackage{fontspec}
\usepackage{amsmath}
\usepackage{amsfonts}
\usepackage{multirow}
\usepackage{array}
\usepackage{mdwlist}
\usepackage{subfig}
\usepackage{floatrow}
%\usepackage{float}
\usepackage{verbatim}
\usepackage{color}
\usepackage{graphicx}
\usepackage{xunicode}
\usepackage{xltxtra}
\usepackage{url}
%\usepackage{dsfont}
%\usepackage{microtype}
\usepackage[htt]{hyphenat}
\usepackage{multicol}
\usepackage{wrapfig}
\usepackage{lipsum}

% FONT SETTINGS ===============================================================

%\setromanfont[Mapping=tex-text]{CMU Serif}
%%\setromanfont[Mapping=tex-text]{CMU Sans Serif} % temporary change until printing
%%\setsansfont[Mapping=tex-text]{CMU Sans Serif}
%%\setmonofont[Mapping=tex-text]{CMU Typewriter Text}
%\setmainfont[Mapping=tex-text]{CMU Serif}
%%\setmainfont[Mapping=tex-text]{CMU Sans Serif}  % temporary change until printing

%\setromanfont[Mapping=tex-text,ExternalLocation=fonts/]{cmunrm.otf}
%\setsansfont[Mapping=tex-text,ExternalLocation=fonts/]{cmunss.otf}
%\setmonofont[Mapping=tex-text,ExternalLocation=fonts/]{cmuntt.otf}
%\setmainfont[Mapping=tex-text, ExternalLocation=fonts/]{cmunss.otf}

\setromanfont{Times New Roman}

% CUSTOM COMMANDS ==============================================================

% Layout macros
\newcommand\spa[1]{\; #1 \;}

% Font macros
\newcommand\resfont[1]{\ensuremath{\mathtt{#1}}}

% Mathematical macros
\newcommand\setmap[3]{#1\{#2 \mapsto #3\}}
\newcommand\getmap[3]{(#2 \mapsto #3) \in #1}
\newcommand\tuple[2]{\ensuremath{\langle#1, #2\rangle}}
\newcommand\mfrac[2]{\ensuremath{\dfrac{#1}{#2}}}
\newcommand\nequiv[2]{\ensuremath{#1 \not\equiv #2}}

% Core Ruby Operational Semantics letter bindings
\newcommand\mem{\mu}

% Core Ruby Operational Semantics low level macros
\newcommand\state[2]{(#1, #2)}
\newcommand\transition[1]{\ensuremath{\overset{#1, c*}{\rightarrow}}}
\newcommand\range[2]{#1, ..., #2}
\newcommand\midrange[5]{\range{#1}{#2}, #3, \range{#4}{#5}}

% Core Ruby Operational Semantics high level macros
\newcommand\operation[5]{\ensuremath{\state{#1}{#2} \transition{#3} \state{#4}{#5}}}
\newcommand\propagation[2]{\operation{#1}{\mem}{#2}{#1'}{\mem'}}

% Core Ruby specific Operational Semantics macros
\newcommand\semicolon[2]{#1; \; #2}
\newcommand\assign[2]{#1 = #2}
\newcommand\mcall[3]{#1.\texttt{#2}(#3)}
\newcommand\ifte[3]{\resfont{if} \; #1 \; \resfont{then} \; #2 \; \resfont{else} \; #3}
\newcommand\newclass[2]{#1.\resfont{new}(#2)}
\newcommand\methoddef[3]{\resfont{def} \; #1(#2) = #3}
\newcommand\classdef[2]{\resfont{class} \; #1 = #2}
\newcommand\with[3]{with \; \tuple{#1}{#2} \; do \; #3}

% Success Typing macros
\newcommand\ssub{\sqsubseteq_S}

% Core Ruby Success Typing letter bindings
\newcommand\classlist{\Delta}
\newcommand\envir{\Gamma}
\newcommand\fields{\Phi}
\newcommand\currclass{l}

% Core Ruby Success Typing inferencing macros
\newcommand\stinfer[5]{\classlist; \; #1; \; \fields \; \underset{\currclass}{\vdash} \; #2: #3 \; \& \; #4; \; #5}


%%%%%%%%%%%%%%%%%%%%%%%%%%% HASKELL STUFF %%%%%%%%%%%%%%%%%%%%%%%%%%%

%\newcommand\typerep[1]{\ensuremath{#1}}
\newcommand\typerep[1]{\lstinline[basicstyle=\normalsize\ttfamily,keywords={}]|#1|}
\newcommand\typefootrep[1]{\textbf{\lstinline[basicstyle=\footnotesize\ttfamily,keywords={}]|#1|}}
% \newcommand\ttyperep[1]{\typerep{#1}}
% \newcommand\mtyperep[1]{\mbox{\typerep{#1}}}

% Arrow types
\newcommand\typeto[2]{\typerep{#1} \typerep{->} \typerep{#2}}
\newcommand\typetotwo[3]{\ensuremath{\typerep{#1} \typerep{->}
                                     \typerep{#2} \typerep{->}
                                     \typerep{#3}}}

\newcommand\tyconapone[2]{\ensuremath{\mbox{\typerep{#1}} \:\: \mbox{#2}}}
\newcommand\tyconaponeC[2]{\ensuremath{\mbox{\typerep{#1}} \:\: \mbox{\typerep{#2}}}}
% \newcommand\tyconapone[2]{\typerep{#1} $\:$ \typerep{#2}}

\newcommand\tyconaptwo[3]{\ensuremath{\mbox{\typerep{#1}} \:\: \mbox{#2} \:\: \mbox{#3}}}


% FLOAT SETUP ==================================================================

% SPELLING =====================================================================

% CODE HIGHLIGHTING ============================================================

\usepackage{listings}

% Define colors that will be used
\definecolor{gray}{rgb}{0.5,0.5,0.5}
\definecolor{darkgreen}{rgb}{0.0,0.5,0.0}
\definecolor{mygreen}{rgb}{0,0.6,0}
\definecolor{mygray}{rgb}{0.5,0.5,0.5}
\definecolor{mymauve}{rgb}{0.58,0,0.82}

% Define common settings for code listings

\lstset{
	backgroundcolor=\color{white},
	basicstyle=\footnotesize,       % style for code
	breakatwhitespace=false,        % sets if automatic breaks should only
									% happen at whitespace
	breaklines=true,                % sets automatic line breaking
	captionpos=b,                   % sets the caption-position to bottom
	commentstyle=\color{mygreen},   % style for comments
	escapeinside={\%*}{*)},         % if you want to add LaTeX within your code
	frame=single,                   % adds a frame around the code
	keepspaces=true,                % keeps spaces in text, useful for keeping 
	keywordstyle=\color{blue},      % keyword style
	numbers=left,
	numbersep=5pt,                  % how far the line-numbers are from the 
									% code
	numberstyle=\tiny\color{mygray},% style for line-numbers
	rulecolor=\color{black},
	stepnumber=1,                   % the step between two line-numbers. If 
									% it's 1, each line will be numbered
	stringstyle=\color{mymauve},    % style for strings
	tabsize=2,                      % sets default tabsize to 2 spaces
}

% Define specific rules for each language

\lstdefinestyle{c}
{
	language=c,
	tabsize=4
}

\lstdefinestyle{haskell}
{
	language=Haskell
}

\lstdefinestyle{ruby}
{
	language=Ruby
}

\lstdefinestyle{erlang}
{
	language=Erlang.
	captionpos=
}

\lstdefinestyle{plain}
{
	captionpos=,
	stepnumber=0
}

% Create new commands for simpler usage

\newcommand\ccode[2]{
	\lstinputlisting[style=c, caption={#1}, label=lst:#2]{src/#2}
}

\newcommand\haskellcode[3]{
	\lstinputlisting[style=haskell, caption={#1}, label=lst:#2]{src/#3}
}

\newcommand\rubycode[2]{
	\lstinputlisting[style=ruby, caption={#1}, label=lst:#1]{src/#2}
}

\newcommand\erlangcode[2]{
	\lstinputlisting[style=erlang, caption={#1}, label=lst:#1]{src/#1}
}

\newcommand\plaintext[2]{
	\lstinputlisting[style=plain, caption={#1}, label=lst:#1]{src/#1}
}

% CHANGE MATH FONT ============================================================

% HYPERREF MUST BE LAST =======================================================

\usepackage[xetex,colorlinks=true,linkcolor=blue,citecolor=darkgreen]{hyperref}

% DOCUMENT INFORMATION =========================================================

\title{Thesis subject} % ===============> FIXME
\author{Όνομα Φοιτητή}
\authoren{Student name}
\datedefense{9}{9}{9999}      % ===============> FIXME
\TRnumber{CSD-SW-TR-*-*} % ok ==================> FIXME
\supervisor{Υπέυθυνος Διπλωματικής}
\supervisorpos{Τίτλος Υπευθύνου}
\committeeone{Πρώτο μέλος επιτροπής}
\committeeonepos{Τίτλος μέλους}
\committeetwo{Δεύτερο μέλος επιτροπής}
\committeetwopos{Τίτλος μέλους}
\committeethree{Τρίτο μέλος επιτροπής}
\committeethreepos{Τίτλος μέλους}

\hypersetup{
	pdftitle={Thesis subject}, % ===> FIXME
	pdfauthor={Όνομα Φοιτητή},
	pdfsubject={},
	pdfkeywords={}
}


% MAIN DOCUMENT ================================================================

\begin{document}

\frontmatter
\maketitle

\def\templen{\parindent}
\setlength{\parindent}{0pt}
\setlength{\parskip}{1.5ex plus 0.5ex minus 0.2ex}
\begin{abstractgr}
	Η υπηρεσία αποθήκευσης δεδομένων μιας υποδομής υπολογιστικού νέφους 
	(cloud computing) έχει κατά κανόνα υψηλές απαιτήσεις σε επιδόσεις, 
	ειδικά όταν είναι επιφορτισμένη με την παροχή των εικονικών δίσκων για 
	τις Εικονικές Μηχανές (VMs) των χρηστών. Σε μεγάλες εγκαταστάσεις όμως, 
	παρατηρείται ότι η επίδοση των σκληρών δίσκων (hard disks) συνήθως δεν 
	επαρκεί για την εξασφάλιση γρήγορων ταχυτήτων για τις ανάγκες των 
	χρηστών. Αυτό οδηγεί τους προγραμματιστές στο να καταφεύγουν σε λύσεις 
	προσωρινής αποθήκευσης δεδομένων (caching) σε γρηγορότερα αποθηκευτικά 
	μέσα, για να βελτιώσουν τις επιδόσεις τους.
	
	Η παρούσα διπλωματική εργασία παρουσιάζει τον \textit{cached}, έναν 
	caching μηχανισμό ο οποίος δημιουργήθηκε για τη βελτίωση της επίδοσης 
	του Archipelago, ενός κατανεμημένου συστήματος αποθήκευσης σε 
	περιβάλλον υπολογιστικού νέφους (cloud computing environment).	
	Επιπροσθέτως, παρουσιάζεται το συμπληρωματικό δικτυακό εργαλείο 
	\textit{synapsed}, το οποίο ανοίγει το δρόμο για την αναβάθμιση του 
	\textit{cached} σε έναν πλήρως κατανεμημένο caching μηχανισμό.
	
	Τα αποτελέσματα από τις πρώτες μετρήσεις είναι αρκετά ενθαρρυντικά και 
	δείχνουν ότι η χρήση του \textit{cached} μπορεί να αυξήσει την επίδοση του 
	Archipelago μέχρι και 400\%.
	\begin{keywordsgr}
		cached, synapsed, synnefo, okeanos, rados, archipelago,  
		εικονική μηχανή, προσωρινή αποθηκευτική μνήμη, Ο(1) 
		πολυπλοκότητα, φυσική μνήμη, κατανεμημένο σύστημα, υπολογιστικό 
		νέφος, κρυφή μνήμη
	\end{keywordsgr}
\end{abstractgr}

\begin{abstracten}
	The performance of the storage service of a cloud infrastructure is very 
	critical, especially when it provides virtual volumes for the Virtual
	Machines of users. Underlying hard disks however, are usually inadequate in 
	providing fast performance for large deployments and the storage engineers 
	commonly have to resort in using various buffering and caching techniques.
	
	This diploma thesis presents \textit{cached}, a caching mechanism that 
	aims to improve the performance of Archipelago, a software-defined, 
	distributed storage layer for cloud computing environments.  
	Additionally, this diploma thesis introduces \textit{synapsed}, a 
	complementary network tool that paves the way for the creation of a 
	fully distributed cache.
	
	Early performance evaluations look very promising and show that 
	\textit{cached} can improve the current Archipelago performance up to 
	400\%.
	\begin{keywordsen}
		cached, synapsed, synnefo, okeanos, rados, archipelago, cloud, 
		virtual machine, volume service, cache, hash table,  O(1) 
		complexity, ram, distributed system, replication, 
		software-defined, RAM
	\end{keywordsen}
\end{abstracten}

\begin{acknowledgementsgr}
	Οι παρακάτω παράγραφοι ήθελα να είναι ένας επίλογος μιας προσπάθειας 
	που ξεκίνησε όταν μπήκα στη σχολή. Όσο όμως τις γράφω και τις σβήνω, 
	νιώθω πως υπάρχουν ακόμα κεφάλαια ανοικτά, που κάνουν τις αναφορές μου 
	σε φίλους και συγγενείς να φαίνονται πρώιμες. Ίσως σε κάποια άλλη 
	σελίδα άλλων ευχαριστιών μπορέσω να μιλήσω για αυτούς τους κοντινούς 
	μου ανθρώπους.

	Το κεφάλαιο που σίγουρα τέλειωσε όμως είναι αυτό της διπλωματικής. Γι' 
	αυτό λοιπόν, θα ήθελα να ευχαριστήσω κάποιους ανθρώπους που συνετέλεσαν 
	καταλυτικά στην ολοκλήρωσή της.
	
	Θα ήθελα αρχικά να ευχαριστήσω τον καθηγητή Νεκτάριο Κοζύρη, που μου 
	έδωσε την ευκαιρία να εργαστώ πάνω σε ένα πολύ ενδιαφέρον και φρέσκο 
	αντικείμενο, το cloud computing.

	Επίσης, θα ήθελα να ευχαριστήσω τον Υ.Δ. Φίλιππο Γιαννάκο για τις φορές 
	που με προσγείωνε όταν χρειαζόταν, που με διόρθωνε όπου είχα λάθος και 
	που μου	έδινε λύσεις όταν βρισκόμουν σε αδιέξοδο.

	Τέλος, θα ήθελα να ευχαριστήσω τον Δρα. Βαγγέλη Κούκη που, παρά τα 
	όποια κενά στις γνώσεις μου, με στήριξε στην επιλογή μου να ασχοληθώ με 
	κάτι τεχνικά δύσκολο και με βοήθησε να τα καλύψω στην πορεία.
\end{acknowledgementsgr}


%\selectlanguage{english}   % Everything from here on is in english
%\setlanguage{english}
%\nogreekalph

\setlength{\parindent}{\templen}
\setlength{\parskip}{0pt}
\tableofcontents
\listoffigures
%\listoftables
\renewcommand{\lstlistlistingname}{List of Listings}
%\lstlistoflistings % changed the title above

\mainmatter
% moved these two commands here so that they don't influence the toc
\setlength{\parindent}{0pt}
\setlength{\parskip}{1.5ex plus 0.5ex minus 0.2ex}

\renewcommand\floatpagefraction{.7}

\begin{frame}[t,plain]
\titlepage
	\note[item]{Καλημέρα σας, ονομάζομαι Αλέξιος Πυργιώτης\\
		Θα σας παρουσιάσω τη διπλωματική μου με τίτλο:..}
	\note[item]{Ακούγεται κάπως περίεργο στα ελληνικά...\\
	αυτό που πραγματεύται είναι την δημιουργία ενός caching μηχανισμού για 
	το Archipelago, ένα distibuted, storage layer}
	\note[item]{Συγκεκριμένα, στην παρουσίαση αυτή θα μιλήσουμε για τον 
		cached, δηλαδή τον caching μηχανισμό μας και αντικείμενο της 
		διπλωματικης, αλλά και για το synapsed, ένα συμπληρωματικό 
		εργαλείο που στόχος του είναι να δώσει στον cached δικτυακές 
		δυνατότητες}
	\note[item]{Σημείωση: Για οικονομία του λόγου, δε θα προβώ σε εξήγηση 
		βασικών όρων όπως VMs, storage. Παρ'όλα αυτά όμως, αν κάποια 
		από αυτές τις έννοιες είναι άγνωστες ή για όποια απορία κατά τη 
		διάρκεια της παρουσίασης, μπορείτε να με διακόψετε και να 
		ρωτήσετε.}
	
\end{frame}

\begin{frame}[t]{Contents}

	\tableofcontents

	\note[item]{O κορμός της παρουσίασης είναι ο εξής:
		\begin{itemize}
			\item Αρχικά, παρουσιάζουμε κάποια εισαγωγικά που 
				αφορούν το background της εργασίας μας.  
				Αναφέρουμε τι είναι το Synnefo και τι είναι η 
				υπηρεσια ~okeanos
			\item Έπειτα, μιλαμε για τον τρόπο που διαχειριζόμαστε 
				το storage και κατ'επέκταση για το archipelago 
				και τον στόχο της διπλωματικής.
			\item Στη συνέχεια μιλάμε για το τι είναι caching και 
				αναφέρουμε κάποιες σύγχρονες λύσεις για 
				caching.
			\item Τα επόμενα δυο κεφάλαια έχουν να κάνουν με τον 
				cached, και συγκεκριμένα με την παρουσίαση της 
				σχεδίασής του και της απόδοσής του
			\item Αντίστοιχα παρουσιάζουμε το synapsed, τη σχεδίαση 
				και υλοποίησή του
			\item Τέλος, συνοψίζουμε όσα ειπώθηκαν παραπάνω και 
				μιλάμε για μελλοντικές εργασίες
		\end{itemize}
	}

\end{frame}

\section{Introduction}

\begin{frame}{Synnefo}
	\note{Ας ξεκινήσουμε με την παρούσα κατάσταση.\\
	Το software που τα ξεκίνησε όλα είναι το Synnefo}

	\includegraphics{images/synnefo-logo.png}

	Open source, production-ready, cloud software.\\
	Designed since 2010 by GRNET.
	\note{\dspc..by GRNET -> Και φυσικά τα παιδιά που βλεπετε εδώ}
	
	\spc
	Synnefo, as most cloud software, has the following services:
	\begin{itemize}
		\item Compute Service
		\item Network Service
		\item Storage Service
		\item Image Service
		\item Identity Service
	\end{itemize}

	\note{\dspc
		\begin{itemize}
			\item Compute service, είναι η υπηρεσία η οποία 
				προμηθεύει τους χρήστες με VMs και επιτρέπει το 
				χειρισμό τους
			\item Network service, είναι η υπηρεσία η οποία δίνει 
				τη δυνατότητα στους χρήστες να δημιουργήσουν 
				ιδιωτικά δίκτυα και να συνδέσουν τα VMs τους σε 
				αυτά.
			\item Storage service, που παρέχει αποθηκευτικό χώρο 
				στους χρήστες.
			 \item Image Service, υπεύθυνο για το deployment ενός 
				 VM από ένα image.  Επίσης, κάνει και 
				 παραμετροποιήσεις (παράδειγμα ssh κλειδιά)
		\end{itemize}
	}
\end{frame}

\begin{frame}{okeanos}

	\includegraphics{images/okeanos-logo.png}

	\begin{itemize}
		\item IaaS service
		\item Targeted at the Greek Academic and Research Community
		\item Designed by GRNET
		\item In production since 2011
	\pause
		\item ...and of course powered by Synnefo.
	\end{itemize}

	\note{
		\begin{itemize}
			\item IaaS είναι πρακτικά η παροχή εικονικής υποδομής 
				σε χρήστες (δηλαδή πάρε υπολογιστή (VM), 
				δίκτυα, αρχεία κτλ)
			\item Δωρεάν για τους Ακαδημαϊκους σκοπούς, ήδη 
				γίνονται εργαστήρια στο EMP και απ' αυτό το 
				εξάμηνο σε άλλες σχολές
			\item \click
			\item Και φυσικά παίζει πάνω σε Synnefo...
		\end{itemize}
	}
				
\end{frame}



\chapter{Chapter 2}\label{ch:ch2}

\section{Section 1}

This section has an important citation\cite{inproceeding}

\subsection{Subsection 1}

This subsection has code in Haskell:
\haskellcode{Sample code}{sample}{samplecode.hs}

It also has a list:
\begin{description}
\item[Item 1]
	First item
\item[Item 2]
	Second item and a footnote\footnote[1]
	{Footnote description.}.
\item[Item 3]
	Third item and text in \textit{italics}.
\end{description}

And an enumerated list:

\begin{enumerate}
\item	First item.
\item	Second item and text in \textbf{bold}
\end{enumerate}

\section{Section 2}

\subsection{Subsection 1}
This subsection has a link to the block of code~\ref{lst:sample} in Section 1.

\subsection{Subsection 2}

% FIXME: Must add content here
This subsection has a FIXME comment, visible only to the author.

\chapter{Archipelago}\label{ch:archip}

This chapter presents Archipelago, a distributed storage layer that is part of 
the Synnefo cloud software. More specifically, Section 
\ref{sec:overview-archip} introduces Archipelago and its basic features.  
Section \ref{sec:arch-archip} shows the architecture of Archipelago and 
explains the components that it consists of. Section \ref{sec:req-archip} 
illustrates how Archipelago handles the I/O request flow up to the storage 
backend and finally, Section \ref{sec:rados-archip} shows an important 
Archipelago storage backend, the RADOS object storage.

\section{Overview}\label{sec:overview-archip}

Archipelago is a distributed storage layer that is part of the Synnefo cloud 
software. It is responsible for creating Copy-on-Write, snapshottable volumes 
for VMs.  Archipelago can be considered as a storage layer (see Figure 
\ref{fig:archipelago_overview_a.pdf}) that is positioned between the VM's block 
device, and the underlying storage. Also, we can see that it is distributed in 
nature since it can run in an arbitrary number of nodes and most importantly 
with zero communication overhead between them.

\diagram{Archipelago overview}{archipelago_overview_a.pdf}

Archipelago has the following objectives:

\begin{itemize}
	\item Thinly provision volumes to VMs with zero data movement.
	\item Snapshot VM volumes and use them as system images with, again, zero 
		data movement.
	\item Allow VM migrations between Archipelago nodes with no restrictions.
	\item Be agnostic to the actual storage backend used.
\end{itemize}

\section{Architecture}\label{sec:arch-archip}

Archipelago has a modular architecture, which allows it to categorize its 
operations and assign them to distinct components. The IPC between these 
components, which are called \textbf{peers} in the Archipelago dialect, is 
facilitated by XSEG.

XSEG can be considered as the kernel of Archipelago and is a custom mechanism 
that:
\begin{inparaenum}[i)]
\item defines a common communication protocol for all peers, regardless of 
	their type (userspace/kernelspace, singlethreaded/multithreaded) and
\item builds a shared memory segment, where peers can share data using 
	zero-copy techniques.
\end{inparaenum}
The above are provided to the peers by the \texttt{libxseg} library.

In Figure \ref{fig:new_sxima_numbered.pdf}, we present the architecture of 
Archipelago.  Moreover, we show the Archipelago peers, the communication 
channels between them and we briefly explain the operations that they are 
responsible for.

\diagram{Archipelago components}{new_sxima_numbered.pdf}

\begin{description}
	\item[XSEG Block Device (xsegbd)] \hfill \\
		xsegbd is a kernel module that exposes a VM's volume as a block 
		device. For each VM, Archipelago registers an xsegbd block 
		device. This block device provides the entry point for the 
		requests that enter the Archipelago layer.
	\item[VoLuMe Composer Daemon (vlmcd)] \hfill \\
		vlmcd accepts requests from the various xsegbd devices and 
		translates them to object requests, with the help of mapperd.
	\item[Mapper Daemon (mapperd)] \hfill \\
		mapperd is responsible for the mapping of volumes to objects.  
		This means that it must tackle a broad set of tasks such as 
		knowing the objects that a volume consists of, cloning and 
		snapshotting volumes and creating new ones.
	\item[Blocker Daemon (blockerd)] \hfill \\
		blockerd is not a specific entity but a family of drivers, each 
		of which is written for a specific storage type. Blockers have 
		a single purpose, to read/write objects from/to the storage.  
		Currently, there are blockers for NFS and the RADOS object 
		storage.
\end{description}

\section{Requests}\label{sec:req-archip}

In the previous section, we have shown that Archipelago uses distinct 
components that operate in various modes and spaces. At this point, one might 
wonder, why can't the above logic be handled solely by our custom block device 
driver?

The reason is because the approach that we have favored gives us the following 
benefits:

\begin{enumerate}
	\item We can do most of our operations in user space instead of kernel 
		space.
	\item We can support multiple storage backends by means of plugging.
	\item We have a more manageable, modular software due to the fact that 
		it is split in smaller, single-purpose entities.
\end{enumerate}

The request flow can be seen in Figure \ref{fig:new_sxima_numbered.pdf}. We 
will explain what happens in every step of the process, with the same order as 
they are presented in Figure \ref{fig:new_sxima_numbered.pdf}:

\begin{enumerate}
	\item The VM issues an I/O request, which is essentially a request to a 
		block device. The hypervisor forwards this request to an xsegbd 
		device, which is attached to every VM.
	\item The xsegbd's purpose is to bring this request to the user space, 
		by copying its data to the shared memory segment and sending a
		request through XSEG to vlmcd
	\item Once the vlmcd receives the request, it forwards it to
		mapperd, along with the volume name. Then, it waits until
		mapperd returns with the objects that this request consists of.  
	\item The mapperd usually has the mapping of a volume, i.e. the objects 
		that a volume consists of, cached in its memory, so it can 
		respond quickly.  Using the volume's mapping, mapperd can  
		decide which are the objects that this request consists of.
		In the rare event that the mapping is not cached, it queries a 
		special purpose blocker to retrieve it.
	\item At this point, mapperd has two courses of action:
		\begin{enumerate}[i)]
			\item it can either respond to vlmcd the correct 
				objects	that it must read from or write to (3) 
				or
			\item perform a CoW operation (5)->(4)->(3)
		\end{enumerate}

		The CoW operation will be performed if the VM has asked to 
		write on an object which shares data with another object.  In 
		this case, mapperd will send a copy request to the data blocker 
		and will create a new object with the same data (5).  Then, 
		mapperd will update its mapping and store it using the blocker 
		for mappings (4). Once all of the above have finished, mapperd 
		will respond to vlmcd the name of the copied object (3).
	\item At this point, vlmcd can send the request for the correct object 
		to the data blocker.
	\item Finally, the blocker will read/write the requested data to the 
		storage type that it has been written for.
\end{enumerate}

Once the storage backend completes the request, the data blocker notifies the 
vlmcd, which in turn notifies the xsebd, which finally notifies the VM.

\subsection{Request polling}\label{sec:arch-poll}

There is a part that has been omitted in the previous section and that is how 
peers send and receive requests. This is a rather intricate aspect of the 
Archipelago IPC, but we will explain it here since it is needed to understand 
mainly the synapsed implementation.

To begin with, one of the features of the shared segment is that it provides a 
number of \textbf{custom ports}. These ports are similar to network ports, in 
the sense that peers bind them and can listen for requests through them. Each 
port has essentially two queues:
\begin{inparaenum}[i)]
\item a \textbf{request queue}, where new requests are accepted, and
\item a \textbf{reply queue}, where replies to requests are received.
\end{inparaenum}

Thus, when a peer wants to send a request to another peer, it merely needs to 
enqueue that request to the request queue of the peer's port. Likewise, when a 
peer responds to a request that was sent by another peer, it will enqueue that 
request in the reply queue of that peer.

Moreover, peers are notified about new requests via signals. When a peer 
receives a signal, it wakes up and then accepts the request. However, context 
switching is expensive, so once the peer has served a request, it will not 
sleep immediately. Instead, it will iterate its ports for new requests for a 
small number of cycles and if no request is received until then, it will sleep.

\section{RADOS and sosd}\label{sec:rados-archip}

As we have mentioned above, Archipelago is not restricted to a storage backend 
and can work with any backend for which a blocker can be created. There is one 
backend however that provides many important features that are used to make 
Archipelago more scalable and reliable.


This backend is RADOS\cite{rados}, which is the object store component of the 
Ceph \footnote{ceph.com} system. Ceph is a free distributed object store and 
file system that has been created by Sage Weil for his doctoral dissertation
\cite{weil-thesis} and has been supported by his company, Inktank, ever since.

The architecture of Ceph can be seen in Figure \ref{fig:ceph.png}.
\footnote{Picture retrieved from the official website on September 24, 2013.  
	All rights go to the respective owners.}

\diagram{Ceph entry points}{ceph.png}

Ceph utilizes RADOS to achieve the following:

\begin{itemize}
	\item \textit{Replication}, which means that there can be many copies 
		of the same object so that the object is always accessible, 
		even when a node experiences a failure.
	\item \textit{Fault tolerance}, which is achieved by not having a 
		single point of failure. Instead, RADOS uses elected servers 
		called \textbf{monitors}, each of which have mappings of the 
		storage nodes where the objects and their replicas are stored.  
	\item \textit{Self-management}, which is possible since monitors know 
		at any time the status of the storage nodes and, for example, 
		can command to create new object replicas if a node experiences 
		a failure.
	\item \textit{Scalability}, which is aided by the fact that there is no 
		point of failure, which means that adding new nodes 
		theoretically does not add any communication overhead.
\end{itemize}

In a nutshell, RADOS consists of the following components:

\begin{itemize}
	\item \textit{object store daemons}, which are userspace processes that run 
		in the storage backend and are responsible for storing the data.
	\item \textit{monitor daemons}, which are monitoring userspace processes 
		that run in an odd number of servers that form a Paxos part-time 
		parliament\cite{Paxos}. Their main responsibility is holding and 
		reliably updating the mapping of objects to object store daemons, as 
		well as self-healing when an object store daemon or monitor daemon has 
		crashed.
\end{itemize}

The mapping of objects to object store daemons is done indirectly with the help 
of the \textit{CRUSH map} and \textit{placement groups (pgs)}.  Generally 
speaking, placement groups can be considered as logical buckets where more than 
one values can be assigned to.

RADOS uses placement groups in the following manner: On initialization, it is 
configured to have a fixed number of placement groups and a range of values 
assigned to each of them. When it is asked to find an object, the object name 
is hashed and its hash value is used to find in which placement group it 
belongs.  The relationship between placement groups and object store daemons is 
stored in CRUSH maps that each monitor daemon holds. This way, the objects are 
pseudorandomly distributed to the storage backend, which in turn implicitly 
guarantees load balancing. 

There are various entry points for RADOS, as is evident from Figure 
\ref{fig:ceph.png}. Archipelago uses librados for I/O requests and more 
specifically, a blocker called sosd\cite{sosd} has been created by Filippos 
Giannakos to facilitate the communication between RADOS and Archipelago.


\begin{comment}
\section{XSEG design}

The shared memory segment provides, besides a common IPC, the following to the 
peers that join it:

\begin{enumerate}
	\item A fast way to synchronize and/or have access to each other's data 
		using relative pointers (XSEG pointers) to the start of the 
		segment.
	\item A custom memory allocation mechanism that allocates data from the 
		segment, taking into account the memory fragmentation issues.
	\item Structures such as hash tables and stacks, that can be used and 
		accessed from different processes.
	\item An indication for the memory that is used by Archipelago.
\end{enumerate}

To achieve the above, XSEG utilizes several components that fall under the 
following main categories:

\begin{description}
	\item[Drivers] \hfill \\
		Drivers implement the common communication and memory 
		allocation protocol for every peer type (userspace/kernelspace, 
		single-threaded/multi-threaded).
	\item[Libraries] \hfill \\
		Libraries (or libxseg to be exact) are linked with peers and 
		allow them to use xtypes (see below) as well as useful generic 
		functions, such as logging.
	\item[Xtypes] \hfill \\
		Xtypes are data structures such as hash tables, queues, stacks 
		etc. whose data reside in the segment and can be used 
		concurrently by all peers. In fact, it is a rule in Archipelago 
		that if a structure must be accessed by more that one peers, it 
		must be designed as an XSEG type that will provide zero-copy 
		guarantees and integration with the segment.
	\item[Peers] \hfill \\
		Peers, which have been discussed above, are the components that 
		are responsible for accepting, processing and sending I/O 
		requests.
\end{description}

\subsection{Drivers}

\subsection{Libraries}

\subsection{Xtypes}\label{sec:arch-xtypes}

The rationale behind xtypes is:

\begin{itemize}
	\item Abstraction(?) layers: Creating inner abstractions layers for 
		software is not a new concept but it's very easy to miss, 
		especially when you start small and end up big.
		
		In a nutshell, when writing code for a new software (in our case 
		a peer for Archipelago but this can apply to most software that 
		surpass the 1000 LOC\footnote[1]
		{Lines Of Code}
		mark) it is wrong practice to create from scratch a monolithic 
		implementation with indistinguishable parts. There is a main 
		reason for this:
		
		Monolithic implementations usually derive from lack of code 
		architecture and planning. Although it is feasible for a 
		programmer to create fully-functional code that meets the 
		necessary requirements, albeit with a lot more effort and 
		concentration, this approach will backfire when the programmer 
		needs to add new features. Since there is no explicit code 
		architecture and the fragile inner correlations are between 
		lines of code and not separate entities, stored precariously in 
		the developer's mind, the result will eventually be constant 
		code refactorization.
		
		One might think that new features happen once in a while in the 
		development cycle but that would be wrong.  This happens more 
		often than you might think and is actually the common case in 
		iteration and test-driven development.

		The right practice instead is to...
	\item Re-usability:...
	\item User-space / Kernel-space agnosticity: (I doubt that such a word 
		even exists...)
\end{itemize}	

\subsection{Peers}\label{sec:arch-peer}

Peers are Archipelago components that are responsible for accepting, processing 
and sending of the I/O requests. They are essential for the modular nature of 
Archipelago since each of them can be considered as a separate entity. They do 
their own logging, signal handling and processing.

The main Archipelago peers can be seen in Figure ?. As we can see from this 
figure, peers are processes that are attached to an XSEG segment. In the 
previous chaptr, we have mentioned that XSEG segments facilitate the IPC between 
different Archipelago components by offering a shared space where process can 
read and write to very fast. This however barely scratches the surface of IPC in 
Archipelago. In the following section, we will discuss more in-length the 
details behind Archipelago IPC

\section{Archipelago IPC}\label{sec:arch-ipc}

An important part of the shared segment that was intentionally skipped in the 
previous section is its ports.

When XSEG is being initialized, it reserves space in the segment for a number 
of custom ports.  These ports are similar to network ports, in the sense that 
peers bind them and can listen for requests through them. Each port has 
essentially two queues:
\begin{inparaenum}[i)]
\item a request queue, where new requests are accepted, and
\item a reply queue, where replies to requests are received.
\end{inparaenum}
XSEG requests do not differ drastically from common requests and typically have 
a target, offset, size and a buffer to store the data. Targets can either be 
objects or volumes, depending on the peer context.

Another important aspect of the Archipelago IPC, is the mechanism with which 
peers poll for new requests, which is seen in the following section.

\subsection{Polling}
\label{sec:arch-poll}

The polling process \ref{sec:arch-poll}is the following:

\begin{enumerate}
	\item On peer initialization, the SIGIO signal is blocked and the PID 
		of the peer is masked from other peers.
	\item On the request loop, the peer checks its ports, without blocking
		and for a limited amount of cycles.
	\item On the last cycle, it unmasks its PID and check its ports one 
		last time.
	\item If there are no enqueued requests, it uses sigtimedwait to sleep 
		for 10s. sigtimedwait is a function that allows to wait 
		synchronously for pending signals, which is generally faster 
		than invoking a signal handler.
	\item If the 10s expire, the peer masks again its PID and goes to step 
		2.
\end{enumerate}

To understand what requests are in XSEG context, we present some parts of the 
xseg\_request structure:

\begin{description}
	\item[source port:] The port from
\end{description}

in particular have basically the following structure:

For instance, queues are  allocated in the segment, both queues and request are 
custom-built entities for XSEG
The interesting part of this process is


Fist of all, we must clarify that in Archipelago, IPC is done strictly between 
peers in the \textbf{same} memory segment. The reason is that we have crafted 
our own methods for IPC and the processes that need it must attain to a certain 
architecture, which is the peer architecture.

The entrance point for IPC is the peer port. When a peer is registered in the 
segment, it attaches itself to a port range. Peer ports are completely different 
to common ports (which are these ports?). When a peer wants to send a request to 
another peer, it must first "get" the registered port on the segment. The xseg 
port is a structure that holds the necessary information as to where to send the 
request. Every port has three different queues; reply, request, free queue.

Request queues are typically a stack that can be addressed from different peers 
in the same segment. For this reason, they are designed as xtypes. For speed 
resons, they are pre-allocated to a certain length and re-allocated on-line, if 
there is need

Also, ports are designed to be considered as paths. That is, when a request is 
sent from one port to another...

\section{Request flow example}

We have bench xseg which works like so:

\begin{enumerate}
	\item Get request
	\item Prepare request
	\item Create chunk
	\item Allocate peer request
	\item Set request (xhash)
	\item Submit request
\end{enumerate}

\subsection{Get request}\label{sec:get-req-archip}

Explain here about xq or in xtypes?
\end{comment}


\chapter{Tiered Storage}\label{ch:tiering}

In this chapter we will discuss the challenges of today's data storage and how 
tiered storage helps in mitigating costs and increasing performance. Moreover, 
we present the current solutions for tiered storage and we evaluate if they can 
be used in conjunction with Archipelago.

\section{Theoretical Background or What is Tiered Storage}

From the genesis of the PC in ? every component was pushed to its limits to 
become as fast as possible. Take for example the CPU chip. The CPU clock is 
blabal and speed of light means 10cm. Has the same happened with disk

\subsection{Caching}

In caching, there are usually the following two policies:

\begin{itemize}
	\item Write-through:
		This policy bla bla bla
	\item Write-back:
		This policy blu blu blu
\end{itemize}

\subsubsection{Eviction}

Caching generally means that you project a large address space of a slow medium 
to the smaller address space of a faster medium. That means that not everything 
can be cached as there is no 1:1 mapping. So, when a cache reaches its maximum 
capacity, it must evict one of its entries

And the big question now arises: which entry?

This is a very old and well documented problem that still troubles the research 
community. It was first faced when creating hardware caches (the L1, L2 CPU 
caches we are familiar with). In 1966, Lazlo Belady proved that the best 
strategy is to evict the entry that is going to be used more later on in the 
future\cite{Belady}.  However, the clairvoyance needed for this strategy was a 
little difficult to implement, so we had to resort to one of the following, 
well-known strategies:

% Mention ehcache approach: 
% http://ehcache.org/documentation/apis/cache-eviction-algorithms
%
% Also look at the following papers
%L. Belady, “A Study of Replacement Algorithms for a
%Virtual-Storage Computer,” IBM Systems Journal, vol.5,
%no.2, pp.78-101, 1966.
\begin{itemize}
	\item \textbf{Random:} Simply, a randomly chosen entry is evicted. This 
		strategy, although it seems simplistic at first, is sometimes chosen 
		due to the ease and speed of each. It is preferred in random workloads 
		where getting fast free space for an entry is more important than the 
		entry that will be evicted.
	\item \textbf{FIFO (First-In-First-Out):} The entry that was first inserted 
		will also be the first to evict. This is also a very simplistic 
		approach as well as easy and fast. Interestingly, although it would 
		seem to produce better results than Random eviction, it is rarely used 
		though, since it assumes that cache entries are used only once, which 
		is not common in real-life situations.
	\item \textbf{LRU (Least-Recently-Used)}
	\item \textbf{LFU (Least-Frequently-Used)}
\end{itemize}

Choosing the LRU strategy is usually a no-brainer. Not only does it 
\textit{seem} more optimal than the other algorithms, but it has also been 
proven, using a Bayesian statistic model, that no other algorithm that tracks 
the last K references to an entry can be more optimal.
%Also, check out this paper:
%An optimality proof of the LRU-K page replacement algorithm
%that proves that no algorithm that keeps track of the K most recent references 
%for a page can be more optimal than LRU.



\section{Comparison between existing methodologies}

\subsection{Block-devices}

There have been many block device tiers that have been built


\subsubsection{Bcache}

Bcache has been designed by Kent Overstreet since 2011 and has been included in 
the Linux kernel (3.10) since the May of 2013.

Bcache allows one to use one or more fast mediums as a cache for slower ones.  
Typically, the slow medium is a RAID array of hard disks and the fast medium 
are SSD drives. Bcache has been specifically built for SSDs and has the 
following characteristics:

\begin{enumerate}
	\item The data are written sequentially and in erase block size 
		granularity, in order to avoid the costly read-erase-modify-write 
		cycle.
	\item It takes special care to mitigate wear-leveling by touching equally 
		all SSD cells
	\item It honors TRIM requests and uses them as hints for its garbage 
		collection.
\end{enumerate}

On to the technical part, bcache is a kernel driver that needs intrusive 
changes to the backing device, rendering previous data unreadable. On a 
nutshell, bcache edits the superblock of both the cache and backing devices in 
order to use them. Then, it exposes to the users a virtual block device, which 
can be formatted to any filesystem. This virtual device's request function 
actually points to the cache device and, when the cache device has reached its 
full utilization, it will flush dirty data to the backing device.

The most striking bcache feature is that it uses a custom built B+tree as an 
index, which has the added benefit that dirty data can be coalesced and flushed 
sequentially to the slower spinning medium. This provides a considerable 
performance speed-up for hard disks since. Some of other noteworthy features of 
bcache are the following:

\begin{enumerate}
	\item It can operate in three modes, write-through, write-back and 
		write-around, which can be switched on/off arbitrarily during normal 
		usage or when the fast medium is congested.
	\item It utilizes a journal log of outstanding writes so that the data are 
		safe, even when an unclean shutdown occurs.
	\item It can bypass sequential IO and send it directly to the backing 
		device, since this workload is tailored for spinning disks.
\end{enumerate}

\subsection{Key-value store}

\subsubsection{Memcached}

\subsection{NoSQL Databases}

\subsubsection{Couchbase}

\subsubsection{Redis}

\subsection{Summary}






\chapter{Design of cached}\label{ch:cached-design}

In the previous chapters, we have addressed the need for tiering in terms of
scalability as well as performance. % Which are these chapters? Link them.

We have also evaluated current caching solutions and described why they couldn't 
be used as a cache tier in Archipelago. % Provide link to a table of comparisons

With the results of chapter ? in mind, we can provide some more strict 
requirements that our solution must have:

% List the requirements for our solution
\begin{enumerate}
	\item \textbf{Nativity:} Our solution must be native to Archipelago i.e.  
		not need any translation layers to communicate with it.
		%explain why
	\item \textbf{Pluggability:} Our solution must be able to provide a 
		caching layer between peers that are already in operating mode 
		without restarting Archipelago. Also, it must be removed 
		without disturbing the service.
	\item \textbf{In-memory:} Our solution must cache requests in RAM, 
		since the next fastest tier, SSDs, are already being used in 
		RADOS as a journal.
\end{enumerate}

For the following chapters, we will drop the \textit{"solution"} moniker and we 
will use instead the proper name of our implementation, "cached", which simply 
means \textbf{cache d}aemon).

The following two chapters are the main bulk of this thesis and they present our 
own implementation that aims to fill the above requirements.

% Foretell what the following sections will contain
More specifically, this chapter provides an in-depth description of the design 
of cached. Section \ref{sec:rationale-design} provides the design rationale of 
cached and explains how its design meets the above requirements. Section 
\ref{sec:comp-design} presents the building blocks of cached while Sections 
\ref{sec:xcache-design}, \ref{sec:xworkq-design} and \ref{sec:xwaitq-design} 
provide a detailed explanation of their design. Section \ref{sec:cached-design}  
presents the interaction between cached and its  as well as the unique 
components that cached consists of.  Finally, in Section ?  we illustrate the 
flow of requests for cached.

\section{Design rationale}\label{sec:rationale-design}

One of the first architectural decisions was to implement cached as an 
Archipelago user-space peer (see Section \ref{sec:arch-peer} about Archipelago 
peers). This choice was the most natural one since it provides the smallest 
possible communication overhead with the other Archipelago peers. Also, this 
design decision covers the nativity requirement we posed at the beginning of 
this chapter.

The above design choice has another advantage too; we can register on-line the 
cached peer between the vlmc and blocker and unregister it when we want to.  
This opens up numerous possibilities such as plugging cached for 
QoS\footnote{Quality of Service} reasons when there is a peak in I/O requests.  
This is possible because, as we have mentioned in Section \ref{sec:arch-ipc}, 
XSEG ports can be registered on-line. Thus, during normal operation, the 
administrator can add the cached port to the request path between vlmc and 
blocker, and all requests will seamlessly be intercepted by cached. This 
follows the same principle with bcache, which plugs its own request\_fn() 
function to the virtual device it creates.  Unlike bcache however, cached can 
be plugged on and off at any time.

This also means that the pluggability requirement is also being met.

The next important design decision was what will cached index. Given that it 
will reside between the vlmc and blocker, where the VM's requests have already 
been translated to object requests, the natural choice is to cache objects.  

The decision to cache objects not only is the most natural one, but also is 
closer to the way our storage (RADOS) handles data. To understand the 
importance of it, consider the following:

Like bcache, our implementation must not only cache object requests fast but 
also try to coalesce them so that, when needed, they will be flushed to the 
slower medium in a more sequential fashion. The fact however that the VMs' 
volumes are partitioned into different objects, means that sequential data (in 
volume context) which reside in different objects will probably not be 
sequential in the storage backend too.

Thus, unlike bcache which expects that the backing device is also the physical 
device and coalesces date accordingly, our implementation is limited only in 
coalescing data in the object range (commonly 4MBs). If our implementation was 
caching in block or volume level, it would be unaware of that fact. 

Having decided that cached will cache objects, the next step is to decide
\begin{inparaenum}[i)]
\item on the index mechanism and
\item on what \textbf{exactly} would we index.
\end{inparaenum}

As for what we would index, it would be an overkill to further partition the 
objects and index regions within them. Moreover, this would make sense only if 
the objects where large (e.g. like volumes). So, our index mechanism should 
index object names solely. As for the index mechanism, we have chosen to use a 
very fast in-memory hash-table for this job. This covers the in-memory 
requirement we have set above. Also, this choice is one of the main reasons 
that our implementation is O(1).

Finally, another important decision was whether cached would be a 
multi-threaded peer. We have decided that we will implement it this way and 
then evaluate the performance of the implementation to find out if we are 
benefited by multi-threading or not.

Thus, cached must be able to work with multiple threads which will accept 
requests from cached's request queue and serve them concurrently with the other 
threads. Of course, multi-threading can be very tricky, especially when we are 
dealing with I/O requests and simultaneous accesses to the same object blocks.  
So, in order to achieve a balance between safety and speed, we use a
fine-grained locking scheme in critical sections that can be seen is discussed 
in detail in Section \ref{sec:xworkq-design}.

\section{Cached components}\label{sec:comp-design}

At this point, we must do an intermission before we show the design of cached.  
Specifically, we will show first the design of the cached's components, since 
many cached operations rely on them and the reader needs prior knowledge of 
them to grasp the cached design.

\subsection{Overview}

In this section, we will list the main components that cached relies on. Per 
Archipelago policy, most of these components have been written in the xtypes 
fashion (see Section \ref{arch-xtypes} about xtypes).  

The components of cached can be seen below:
 
% TODO: Link List to the definitions
\begin{itemize}
	\item xcache, an xtype that provides indexing support, amongst many 
		other things
	\item xworkq, an xtype that guarantees atomicity for execution of jobs 
		on the same object
	\item xwaitq, an xtype that allows conditional execution of jobs
\end{itemize}

and their design will be discussed in-depth in the following sections.

Also, we must note that the above components predate our cached implementation 
and are not a contribution of this thesis\footnote{xcache is an exception since 
	we have extended its functionalities for our purposes}. They are 
presented however in this thesis for clarity reasons. 

\subsection{The xcache xtype}\label{sec:xcache-design}

xcache is the most important component of cached. It is responsible for several 
key aspects of caching such as:

\begin{itemize}
	\item entry indexing,
	\item entry eviction,
	\item concurrency control and
	\item event hooks
\end{itemize}

Below, we can see a design graph of xcache:

\fixme add Figure here

% While explaining, point to the diagram objects, a, b, c etc.
\fixme add better design explanation \\
As we can see above, xcache utilizes two hash tables. One hash table is 
responsible for indexing entries (or more generally speaking "cache entries") 
that are active in cache.  The other hash table is responsible for indexing 
evicted cache entries that have pending jobs.  Again, more generally speaking, 
evicted cache entries are entries whose refcount has not dropped to zero yet.

On the following subsections, we present the features of xcache as well as 
their design.

\subsubsection{Entry Preallocation}\label{sec:xcache-entry-design}

Since xcache indexes a bounded number of entries, there is no need to allocate 
them on-the fly using malloc/free. Considering that we are caching at RAM level 
and not at SSD level, the system call overhead will have a considerable impact 
on performance. Thus, in our case, we pre-allocate the necessary space in 
advance and store them in a cache-node pool (note that this is different from 
the bucket pool).

\subsubsection{Entry indexing}\label{sec:xcache-index-design}

The index mechanism that xcache uses is a hash table named xhash, also an 
xtype. The reason why a hash table is used as an index mechanism is because:

\begin{enumerate}
	\item Given that we index only a certain number of entries, we expect 
		the that the insert, lookup and delete operations are in 
		constant time (see below for an explanation why)
	\item Hash tables can preallocate the space needed whereas 
		tries/b-trees/bst will allocate nodes as new entries are 
		inserted. Again, the fact that we index a certain number of 
		entries means that we expect that we will have many evictions 
		and insertions.
	\item We don't need to do substring matches (advantage of tries)
	\item We don't need to traverse the entries sequentially (advantage of 
		B-trees and BSTs)
\end{enumerate}

The hash table that is used is heavily based on dictobject\cite{dictobject},
the Python dictionary implementation in C. Distobject has been created to 
minimize the collisions and the hops (\fixme Explain that better). Its only 
drawback is that it needs to resize when the table's entry history has reached 
the 2/3 of its capacity.

Besides the hash table, which answers to the question "Where is the entry?" we 
also need another mechanism to answer the question "Is the entry still 
referenced?". xcache has such a mechanism which is commonly called "reference 
counting". Specifically, each entry has a counter that is 
incremented/decremented when a user accesses/releases an entry.

To sum up, when an entry is inserted, we use its name as a key and we update 
its refcount to 2, one reference from the user and one standard reference from 
the hash table. When we lookup for an entry, we use the entry's name as a key 
and then increment by 1 its refcount. 

\subsubsection{Entry eviction}\label{xcache-eviction-design}

The decision to have xcache index a bounded number of entries means that when 
it reaches its maximum capacity and is requested to index a new entry, it has 
to resort to the eviction of a previously cached entry. Evicted entries are not 
removed immediately from xcache. They are instead set in an "evicted" state and 
they reside in a special-purpose hash table until the user confirms that they 
can be removed. 

xcache handles evictions in an interesting way. More specifically, evictions 
occur implicitly and not explicitly, meaning that the user (peer) doesn't have 
to evict entries manually. For example, when a user tries to insert a new entry 
to an already full cache, the insertion will succeed and the user will not be 
prompted to evict an entry manually. Moreover, the user will be notified via 
specific event hook that is triggered upon eviction.

The scheme of implicit evictions and later on notification of the user has the 
advantage that lookups, inserts and evictions can occur atomically by xcache.  
This wouldn't be the case if the user was responsible for the evictions.

As for the eviction strategy, we have utilized an LRU queue. Not only it's 
optimal (\fixme verify it) for our purposes, but we have also mitigated the 
cost of keeping the last references for each entry by creating a simple LRU 
algorithm, which has O(1) complexity for all update actions. More about the 
implementation of the LRU algorithm can be found in Section 
\ref{xcache-evict-imp}.

\subsubsection{Concurrency control}

The concept of concurrency control has been discussed in chapter ?. The goal of 
xcache is to handle safely - and preferably fast - simultaneous accesses to the 
shared memory.

In order to do so, we must first identify which are the critical sections of 
xcache, to wit, the sections where a thread modifies a shared structure. These 
sections are the following:

\begin{itemize}
	\item
		\textbf{Most xhash operations:} Inserts and removals can modify 
		the hash table (e.g. they can resize it, add more entries or 
		delete existing ones). This also means that lookups must not 
		run simultaneously with the above two operations.
	\item
		\textbf{Cache node claiming:} Before an entry is inserted, it 
		must	acquire one of the pre-allocated nodes from the 
		cache-node pool and we must ensure that this can happen 
		concurrently from all threads.
	\item
		\textbf{Entry migration:} An entry can migrate from one hash 
		table to the other e.g. on cache eviction. This migration 
		involves a series of xhash operations; removal from one hash 
		table and subsequent insertion to the other. These two 
		operations must occur atomically.
	\item
		\textbf{Reference counting:} Every entry must have a reference 
		counter.  Reference counters provide a simple way to determine 
		when an entry can be safely removed. Since many threads can 
		have access to the same entry, we must provide a way to update 
		the reference counters atomically.
	\item
		\textbf{LRU updates:} Most actions that involve cache entries 
		must subsequently update the LRU queue. The updates at the LRU 
		queue must also occur atomically.
\end{itemize}

Let's see what guarantees we provide for each of the above scenarios:

\begin{itemize}
	\item
		\textbf{xhash operations:} We provide a lock for each hash 
		table. Only one thread can access each hash table at any time.
	\item
		\textbf{Cache node claiming:} The cache-node queue, is also 
		protected by a lock.
	\item
		\textbf{Entry migration:} When an entry is migrated from one 
		hash table to the other, we always acquire the lock of the hash 
		table of active entries and then the lock of the hash table of 
		the evicted entries. The order on which we take the locks is 
		very strict to avoid deadlocks.
	\item
		\textbf{Reference counting:} For the atomic increases and 
		decreases of a counter, we don't need a lock and its added 
		overhead. Instead, we can use the atomic get and atomic put 
		operations that the CPU provides.
	\item
		\textbf{LRU updates:} Since the majority of LRU updates take
		place when a new entry is inserted in the hash table, we can 
		protect our LRU under the same cache lock.
\end{itemize}

\subsubsection{Re-insertion}

We have previously mentioned that in xcache, there can be data migration 
between hash tables. Most commonly, an entry that is evicted from the active 
cache entries migrates to the hash table of the evicted cache entries, until 
its reference count falls to zero and can be freed.

However, what happens when xcache receives a request for an evicted entry which 
hasn't been freed yet? 

In this case, the entry switches state again and is inserted back to the hash 
table of active entries. Also, its reference counter is incremented accordingly 
in order not to be freed amidst this process.

This way, we can avoid waiting for an entry that has just been evicted to flush 
its data. 

\subsubsection{Event hooks}\label{sec:xcache-hooks-design}

Since xcache is created to provide core caching functionalities for other 
peers, it must also notify them when it takes an implicit action that the peer 
is not aware of. In Section \ref{xcache-eviction-design} we have seen one 
implicit action that xccahe takes when a user inserts an entry, namely 
eviction. 

Besides this event, there are others. The complete list is the following:

\begin{description}
\item[cache node initialization:]
	This hook is triggered when a cache node is initialized. It is 
	triggered once only for each node, during the initialization phase of 
	xcache.
\item[cache entry initialization:]
	This hook is triggered when a cache entry has been inserted in the 
	cache.
\item[cache entry eviction:]
	This hook is triggered when a cache entry has been evicted from the 
	cache.
\item[cache entry reinsertion:]
	This hook is triggered when an evicted entry has been reinserted in the 
	cache.
\item[cache entry finalization:]
	This hook is triggered when an evicted entry's refcount has dropped to 
	0. This serves as a warning for the user who has the opportunity to let 
	the cache entry go or increment its refcount.
\item[cache entry put:]
	This hook is triggered when an evicted entry has been totally removed 
	from the cache.
\item[cache entry free:]
	This hook is triggered when a removed entry's cache node has been sent 
	back to the cache node pool.
\end{description}

For each of the above events, we have created the respective event hook. The 
peer that uses xcache may choose, if it wants, to use them and if so, it can 
plug its own event function for each hook which will be called when the event 
is triggered.

\subsection{xcache flow}

To make the way xcache works a bit more clearer, we will see the flow for three 
of the main xcache operations; lookup of an entry; insertion of a new entry and 
removal of an entry:

\subsubsection{Insertion}

\fixme add figure and explanation

\subsubsection{Lookup}

\fixme add figure and explanation

\subsubsection{Put}

\fixme add figure and explanation

\subsection{The xworkq xtype}\label{sec:xworkq-design}

The xworkq xtype is a useful abstraction for concurrency control. Its purpose 
is to enqueue "jobs" (protected by a lock) and ensure that only one thread will 
execute them. There is no distinction as to which thread this will be, as well 
as no execution condition. The executive thread is simply the one that acquires 
the lock first.

xworkq is generally used when multiple threads want simultaneous access to a 
critical section. Instead of spinning indefinitely, waiting for a thread to 
finish, they can enqueue their job in the xworkq and resume processing other 
requests. xworkq is also generic by nature, since the "job" is simply a target 
function and its input data.

On the following figure, we can see the design of xworkq:

% Design of xworkq
\fixme add figure

It consists of a queue where jobs are enqueued.  The thread that enqueues a job 
can attempt to execute it too, by acquiring a lock for the xworkq. If the lock 
is unheld, the thread will acquire and will be able to execute the enqueued 
job. Else, it can safely leave and its job will be executed by the thread that 
has holds the lock.

In cached context, every object has an xworkq. Whenever a new request is 
accepted/received for an object, it is enqueued in the xworkq and we are thus 
ensured that only one thread at a time can have access to the object's data and 
metadata.

\subsection{The xwaitq xtype}\label{sec:xwaitq-design}

The xwaitq xtype bears some similarities to the xworkq xtype. Like xworkq, it 
is also an abstraction where "jobs" are enqueued and dequeued later on to be 
executed. Unlike xworkq though, jobs are executed only when a predefined 
condition is met. Another distinction is that the jobs in xwaitq are considered 
to be thread-safe and can be executed concurrently by any thread.  

xwaitq is commonly used in non-critical code sections that can be executed only 
under specific, predefined circumstances. The "jobs" that are enqueued in 
xwaitq are the same as the jobs of xworkq.

% Design of xwaitq
\fixme add figure

Unlike xworkq, before a job is enqueued, the thread can attempt to execute it 
by checking the execution condition. Only if the condition is \textbf{not} met 
does the thread enqueue the job to the queue. Before the thread leaves, it 
"signals" the queue and essentially rechecks the condition to ensure that it 
can't be executed. It can then safely leave since its job will be executed when 
another thread signals the queue successfully.

In cached context, xwaitqs are used to enqueue jobs which cannot be executed 
immediately. Common cases are when we have run out of space, when we have run 
out of requests etc.

\section{Cached Design}\label{sec:cached-design}

At this point, we have discussed in length the design of the cached components.  
Having the above sections in mind, we can proceed with presenting how cached 
has been designed.

Cached has been been designed mainly as the orchestrator, a peer that utilizes 
several different components to handle various tasks such as indexing (xcache), 
concurrency (xworkq) and deferred/conditional execution (xwaitq). Cached 
however is not limited to the above role as these components do not cover all 
of the needed tasks. There are several other key tasks that cached must 
undertake, namely:

\begin{enumerate}
	\item Request handling
	\item Write policy enforcing
	\item Cache coherence
\end{enumerate}

Moreover, cached extends its repertoire using some unique components, namely:

\begin{enumerate}
	\item Bucket pool
	\item Cache-IO
	\item Book-keeping
\end{enumerate}

We will illustrate the design of cached from two different perspectives: the 
operational perspective, which can be seen in Figure 
\ref{fig:cached-design.pdf}, and the component perspective, which can be seen 
if Figure ?. Moreover, we will further explain how cached manages the above new 
components and tasks in the following sections.

\diagram{cached-design.pdf}{Cached design}

\fixme add component diagram for cached (xcache, xworkq, objects, buckets etc.)

\subsection{Request handling}

Cached operates as a peer that receives requests from the vlmc. The majority of 
these requests will be read/write requests, but there are other types of 
requests too such as copy requests (sent when an object is copied-on-write) and 
info requests i.e. queries on what is the size of an object. Each of these 
requests must be handled independently, using special-purpose functions.

Furthermore, cached will also issue requests to the blocker mainly on two 
occasions: when it flushes a dirty object and when operating in write-through 
mode.

This means that, cached must be able to categorize requests and send them to 
the appropriate functions. Moreover, it must be able to create requests of its 
own, as well as handle cases such as running out of requests.

\subsection{Write policy enforcing}\label{sec:cached-wcp-design}

The user defines beforehand what will the write policy of cached be. There are 
two options: write-through and write-back. These policies aren't new and have 
already been discussed in chapter ?, but let's see what these policies 
translate to in cached context.

\begin{itemize}
	\item
		In \textbf{write-back} mode, cached caches writes, immediately 
		serves the request back and marks the data as dirty. When a read 
		arrives, it either serves the request with the dirty data 
		(read-hit) or forwards the request to the storage peer and 
		caches the answer (read-miss).

		This policy is used when we want to improve read and write speed 
		and can sacrifice data safety.
	\item
		In \textbf{write-through} mode, cached forwards writes to 
		blocker, servers the request when blocker replies, caches the 
		data and marks them as valid.  When a read arrives, it either 
		serves the request with the valid data (read-hit) or forwards 
		the request to the storage peer and caches the answer 
		(read-miss).

		This policy is used when we want to improve read speed and want 
		to make sure that no data will be lost.
\end{itemize}	

These policies are specified once during cached's deployment and cannot be 
switched on/off later on.

\subsection{Cache coherence}

Cache coherence is a concept that is closely connected with concurrency 
control, as is evident in Figure \ref{fig:cached-design.pdf}. It is commonly 
used as a term in the literature on the subject of caches for multiprocessor 
systems or distributed caches and it refers to:

\begin{enumerate}
	\item the consistency of data that are stored in the cache and
	\item how each change is propagated through the system, to the other 
		storage tiers.\label{list:second-coherence}
\end{enumerate}

If we wanted to rephrase the above to match our case, we could say that it 
refers to the consistency of the data that are stored in cached and how each 
change is propagated to the blocker.

The consistency of the cached data is greatly secured using the xworkq as the 
guard of parallel accesses to the object's data. However, the guarantee that 
data updates are conducted serially does not suffice for the consistency 
requirement.

The problem is evident if you consider the eviction of an object with dirty 
data. xcache does not operate on object level and thus is unaware of the 
contents of the cache entry. Thus, cached must use the cache entry finalization 
hook to increment the refcount of the object so that it can not be removed 
until all data have been flushed to the blocker. Essentially, cached overrides 
the xcache's book-keeping to maintain cache coherency.

The propagation of the data to the blocker is solely cached's job. One might 
consider that flushing the dirty data of an object would be sufficient to 
propagate the changes to the blocker. However, consider the scenario when an 
object is evicted and then subsequently reinserted, overwritten and evicted 
again. The result would be two flush requests, for the same data, sent to the 
blocker. Although the blocker guarantees that there will be no page-tearing, we 
cannot be sure about the order on which the data will be written.

To solve this problem, flush jobs must be deferred (using an xwaitq) if another 
flush for the same object is in flight. 

\subsection{Bucket pool}

% Explain what buckets are:
% a) That they are bounded
% b) That they have states
% c) That they are preallocated and every thread enqueues and dequeues an index 
% to them

We have already explained in Section \ref{sec:rationale-design} the reason why 
cached caches objects. There is, however, one important design issue that we 
must address. The issue is how will cached perceive the object's data. To make 
it a bit clearer, given that an object typically has 4MB of size, what should 
cached do when a user requests e.g. a 16KB chunk of it?

If we perceived the object's data as a monolithic chunk, we would need to read 
and cache the whole object just to reply to the user's request. If the user 
then requests a chunk from another object, we would have to cache that object 
too and in the end, we would thrash our cache
\footnote{
	cache thrashing occurs when we aggressively cache data that is only 
	used once and effectively leads to a snowball of evictions
}.

The solution we propose is to further divide objects to the next and final 
logical entity, \emph{buckets} (typically 4KB of size). Each bucket consists of 
its data and metadata and cannot be half-empty, or half-allocated. This way, we 
can also know which parts of the cached object are actually written, or are in 
the process of being read etc.

Thus, our solution to the hypothetical problem we have posed above is to 
request 16KB from the blocker, store the result in 4 buckets and then respond 
the request to the user.

The above answer, however, is not entirely complete. It implies that buckets 
are something readily available or attached to the object. Although each object 
\textit{could} have its buckets pre-allocated, this would limit the objects 
that we can cache since that, even if the user requested only a small chunk, 
each object would statically need 4MB of space.

Ideally, we would like to be able to cache thousands of objects but
\begin{inparaenum}[i)]
\item allocate a much smaller amount of buckets and
\item strictly when the user requests to.
\end{inparaenum}
To make things even faster, we would also like the buckets to be preallocated 
(like cache nodes in Section \ref{sec:xcache-entry-design}) to avoid the 
overhead of malloc/free system calls.

We have achieved the above by creating a bucket pool. The design of the bucket 
pool is the following:

\fixme add diagram

The bucket pool size is static and has been set during initialization by the 
administrator. After the necessary space has been allocated, it is divided in 
buckets and all the bucket indexes are pushed in a lock-protected stack (xq).  
Then, each thread can pop bucket indexes from the pool and attach them
to an object when needed. When that object is evicted, its attached buckets 
indexes are pushed back to the pool.

\subsection{Cache-IO}

When a request is accepted, all Archipelago peers commonly embed it in a peer 
request (you can see more about peer requests here ?). The peer request always 
holds the original request and optionally, several other fields that are of 
importance to the peer.

In our case, we keep in cached's peer requests a new structure called Cache-IO, 
on which we store the xcache handler of the request handler as well as the 
state of the request. Moreover, since the request may break internally in many 
others, we keep track of how many pending requests remain until the original 
one can be completed.

\subsection{Book-keeping}

In order to know when an object is in flushing state, when a bucket has been 
allocated or when a bucket is being read, cached must employ some sort of 
book-keeping of their states. The entities that require to track their states 
are:

\begin{enumerate}
\item \textbf{buckets:} Each bucket has two different states that must be 
	tracked. The first is its allocation state:
	\begin{enumerate}[i)]
		\item free, meaning that the bucket is not allocated
		\item claimed, meaning that the bucket is allocated
	\end{enumerate}
	The second is its data state:
	\begin{enumerate}[i)]
		\item invalid, meaning that it currently holds no data
		\item valid, meaning that it has data that correspond with the 
			blockers data
		\item loading, meaning that a read request has been sent to the 
			blocker to be filled with data
		\item dirty, meaning that it currently holds newer data that 
			what the blocker has.
	\end{enumerate}
\item \textbf{Objects:} The bucket states of an object provide a good 
	indication of the object's status, yet not a complete one. The statuses 
	we keep for the objects are:
	\begin{enumerate}[i)]
		\item ready, meaning that it is ready to accept data
		\item invalidated, meaning that it is not ready to accept data
		\item flushing, meaning that it is currently flushing dirty 
			data to the blocker
		\item failed, meaning that a request has failed for this object 
			and we must stop using it
	\end{enumerate}
\item \textbf{Cache-ios:} Cache-ios also have states. They are the 
	following:
	\begin{enumerate}[i)]
		\item accpted, meaning that the request has just been accepted
		\item reading, meaning that the request wants to read data
		\item writing, meaning that the request wants to write data
		\item served, meaning that the request has been served the data 
			it needed
		\item failed, meaning that the request has failed
	\end{enumerate}
\end{enumerate}

Moreover, we keep global and per-object counters of every object's bucket 
states. The benefits from this approach is that we can know at any time if an 
object (or cached in general) has a bucket in a certain state.

\section{Cached Operation}\label{sec:op-design}

\subsection{Write}

This is the flow for the write path:

\fixme add diagram and explanation

\subsection{Read}

This is the flow for the read path:

\fixme add diagram and explanation

\subsection{Copy}

This is the flow for the read path:

\fixme add diagram and explanation

\subsection{Info}

This is the flow for the read path:

\fixme add diagram and explanation


\chapter{Implementation of cached}\label{ch:cached-implementation}

In the previous chapter, we presented a design overview for cached and its 
components. In this chapter we will blabla how the above design has been
implemented and explain in depth the structures and functions that have been 
created for this purpose.

More specifically, sections ? - ? provide implementation information for the 
components of cached, as described in Chapter ?. Next, section ? presents the 
actual initialization and blabla operations using excerpts from the code.

\section{Implementation of xcache}

In this section, we describe how we implemented the design concept of section 
\ref{sec:xcache-design}. The main \xcache structure is the following:

\ccode{Main \xcache struct}{xcache-struct.h}

Each of the above \xcache struct fields serves a design purpose.
Let's see which fields help in what:

\subsection{Entry Preallocation}

The relevant fields for this purpose can be seen in the following code listing:

\ccode{\xcache struct fields for preallocated entries}{xcache-prealloc.h}

and the definition of the \texttt{xcache entry} struct which shows up in 
\texttt{xcache} struct can be seen below:

\ccode{\texttt{xcache entry} struct}{xcache-entry.h}

Let's start by listing what \texttt{xcache entry} consists of. First of all, it 
must have a name. Since we preallocate the entries and cannot know in runtime 
their length, we must allocate as much space as possible. The \texttt{char 
	name[XSEG\_MAX\_TARGETLEN + 1]} field, which is 256 characters long, is 
long enough to hold the target's name. Also, as we have mentioned in Section 
\ref{sec:entry-prealloc-design}, xcache must be agnostic of the cache contents.  
To this end, we use the generic \texttt{void *priv} field as a pointer to the 
actual entry content. The rest of the fields will be explained in the following 
chapters.

Let's continue now with the fields of Listing \ref{lst:xcache-prealloc.h}. Since 
we preallocate the entries using \texttt{malloc}, they take up a contiguous 
space in memory.  The start of this space is the where the \texttt{*nodes} field 
points to. The \texttt{free\_nodes} field works similarly to the free\_entries 
field in Section \ref{sec:get-req-archip} i.e. it is a stack where indexes to 
unused nodes are pushed. These indexes will be seen in various code excerpts in 
this chapter and have a specific name, 
\texttt{xcache\_handler}\footnote{\#define xcache\_handler uint64\_t}.

\subsection{Entry Indexing}

The relevant fields for this purpose can be seen in the following code listing:

\ccode{xcache struct fields for entry indexing}{xcache-index.h}

As we have mentioned in Section \ref{sec:xcache-index-design}, we utilize two 
hash tables, one for the cached entries and one for the former cached entries 
(or evicted entries or removed entries). These hash tables can be accessed from 
the \texttt{xcache struct} and are \texttt{*entries} and \texttt{*rm\_entries} 
respectively.

These are the functions which are related to indexing and \xcache exposes to the 
peer function:

\ccode{Indexing functions}{xcache-index.c}

All of these function need a pointer to the \xcache struct. Here's a brief 
description of them:

\begin{description}
	\item[xcache\_lookup:]
		Takes the target's name as an argument and searches for it in 
		cache.\\
		Returns on failure: NoEntry\footnote{\#define NoEntry 
			(xcache\_handler)-1}\\
		Returns on success: the requested handler.\\
		\textbf{Note:} Looks only in \texttt{entries}.
	\item[xcache\_insert:]
		Takes the handler of an allocated entry as an argument and uses 
		it to index that entry.\\
		Returns on failure: NoEntry.
		Returns on success:
		\begin{inparaenum}[\itshape a\upshape)]
		\item the same handler or
		\item another one, if this entry already exists in cache.
		\end{inparaenum}
		\textbf{Note:} It looks up first if the entry exist in 
		\texttt{entries} or \texttt{rm\_entries}. The later case can 
		lead to re-insertions.
	\item[xcache\_remove:]
		Takes the handler of an allocated entry as an argument and uses 
		it to remove that entry.\\
		Returns on failure: -1.
		Returns on success: 0.
		\textbf{Note:} Removes entries only from \texttt{entries} hash 
		table.
	\item[xcache\_invalidate:]
		An \texttt{xcache\_remove} spin-off. Takes the name of the entry 
		as an argument, looks it up and then removes it
		Returns on failure: -1.
		Returns on success: 0.
		\textbf{Note:} Unlike remove, entries can either be on 
		\texttt{entries} or \texttt{rm\_entries} hash table.
\end{description}

\subsection{Entry eviction}

The relevant fields for this purpose can be seen in the following code listing:

\ccode{\xcache struct fields for eviction}{xcache-evict.h}

As we have mentioned in Section \ref{sec:xcache-evict-design}, we resort to 
eviction when the cache is full and new entries can't be inserted. This entry is 
the Least Recently Used entry. The doubly-linked list we maintain for this end 
can be seen below:

\ccode{Doubly-linked LRU list}{xcache-dlist.h}

Lets explain these fields a bit:

\begin{description}
	\item[lru:] Obviously, it's the least recently used entry. It can be 
		considered as the one end of the doubly linked list.
	\item[mru:] The entry that has just been used. It can be considered as 
		the other end of the doubly-linked list
	\item[younger:] This entry-specific field points to an entry used right 
		after our entry was used.
	\item[older:] Same as "younger", it points to the entry that has been 
		used right before our entry was used.
\end{description}

Finally, as we have explained in Section \ref{sec:xcache-evict-design}, the 
eviction internals should normally not bother the user. However, if the user 
wants to, \xcache provides the exposes the following functions:

\begin{description}
	\item[xcache\_evict\_lru:] The name says it all, it evicts the recently 
		used item.
	\item[xcache\_peek\_and\_get\_lru:] This function allows the user to 
		atomically take a peek on the Least Recently Used entry and also 
		update its refcount.
\end{description}

\subsection{Concurrency control}\label{sec:conc-imp}

\paragraph{Locking}\label{par:lock}

\paragraph{Reference counting}\label{par:refcount-imp}

The refcount model in xcache should be familiar to most people:

% Turn this to figure
\begin{itemize}
	\item When an entry is inserted in cache, the cache holds a reference 
		for it (ref = 1).
	\item Whenever a new lookup for this cache entry succeeds, the reference 
		is increased by 1 (ref++)
	\item When the request that has issued the lookup has finished with an 
		entry, the reference is decreased by 1. (ref--)
	\item When a cache entry is evicted by cache, the its ref is decreased 
		by 1. (ref--)
\end{itemize}

Some common refcount cases are:

\begin{itemize}
	\item active entry with pending jobs (ref > 1)
	\item active entry with no pending jobs (ref = 1)
	\item evicted entry with pending jobs (ref > 0)
	\item evicted entry with no pending jobs (ref = 0)
\end{itemize}

\begin{table}[tbp]
	\centering
	\begin{tabular}{ | l | l | }
		\hline
		Case & Refcount \\ \hline \hline
		active entry with pending jobs & ref > 1 \\ \hline
		active entry with no pending jobs & ref = 1 \\ \hline
		evicted entry with pending jobs & ref > 0 \\ \hline
		evicted entry with no pending jobs & ref = 0 \\ \hline
	\end{tabular}
	\caption{Reference counting of xcache}
	\label{tab:refcount}
\end{table}

and, as always, the entry is freed only when its ref = 0.

\subsection{Event hooks}

The hooks that xcache provides to users are:

\begin{itemize}
	\item on\_ init: called on cache entry initialization.
	\item on\_put: called when the last reference to the cache entry is put
	\item on\_evict: called when a cache entry is evicted.
	\item on\_node\_init: called on initial node preparation.
	\item post\_evict: called after an eviction has occurred, with cache  
		lock held.
	\item on\_free: called when a cache entry is freed.
	\item on\_finalize: called to hint the user that the cache entry's ref 
		has dropped to zero.
	\item on\_reinsert: called when a cache entry has been in cache
\end{itemize}


\section{Implementation of cached}\label{sec:cached-imp}

The cached is the following:

\ccode{Main cached struct}{cached.h}

and the cached entries are the following

\ccode{Cahed entry struct}{ce.h}

\section{Bucket pool}

And the bucket implementation is the following:

\ccode{Bucket implementation}{bucket.h}

\subsection{Bucket/Object states}

Every object has a state, which is set atomically by threads. The state list is
the following:

\begin{itemize}
	\item READY: the object is ready to be used
	\item FLUSHING: the object is flushing its dirty buckets
	\item DELETING: there is a delete request that has been sent to the 
		blocker for this object
	\item INVALIDATED: the object has been deleted
	\item FAILED: something went very wrong with this object
\end{itemize}

Also, object buckets have their own states too. These are divided in allocation 
states:

\begin{itemize}
	\item FREE: the bucket has not been touched or has been flushed
	\item CLAIMED: the bucket has been claimed
\end{itemize}

and data states:

\begin{itemize}
	\item INVALID: the same as empty
	\item LOADING: there is a pending read to blocker for this bucket
	\item VALID: the bucket is clean and can be read
	\item DIRTY: the bucket can be read but its contents have not been
		written to the underlying storage
	\item WRITING: there is a pending write to blocker for this bucket
\end{itemize}

Also, for cios we have the following:

\begin{itemize}
	\item FAILED: at least one of the cio request(s) has failed
	\item ACCEPTED: normal mode
	\item READING: at least one of the cio request(s) is pending a read
	\item WRITING: at least one of the cio request(s) is pending a write
	\item SERVED: all requests have been served
\end{itemize}


Finally, for every object there are bucket state counters, which are increased/
decreased when a bucket state is changed. These counters give us an O(1)
glimpse to the bucket states of an object.




\backmatter
\cleardoublepage % start at the next odd page
\phantomsection % correct hyperlinking
\bibliography{references}
\bibliographystyle{softlab-thesis}
% \include{glossary}
% \chapter{Appendix}
% \printindex

\end{document}
